\color{unidarkgreen}{\subsection{Фиксация положения выключателей В1 и В2 (ФПВ В1, ФПВ В2)}\label{func:fpv}}
\color{black}

\begin{enumerate}

\item
Функциональный блок фиксации положения выключателей (ФПВ) предназначен для определения состояния выключателей на основании технологической информации:

\begin{itemize}
    \item состояния блок-контактов выключателя;
    \item положения разъединителей в цепи выключателя;
    \item информации о выводе в ремонт выключателя.
\end{itemize}

%%%%%%%%%%%%%%%%%%%%%%%%%%%%%%%%%%%%%%%% УСЛОВИЯ %%%%%%%%%%%%%%%%%%%%%%%%%%%%%%%%%%%%%%%%%%%%%%%%%%%%%%%%
\def\Qone{Не определено}
\def\Qtwo{Не определено}

\ifbool{isLV}{
\def\Qone{В1} 
\def\Qtwo{В2} 
}{}

\ifbool{isDZAT}{
\def\Qone{В1(ЛВ) ВН} 
\def\Qtwo{В2(ОВ) ВН} 
}{}

\ifbool{isLV}{
\item
Контроль положения разъединителей при формировании сигналов включенного или отключенного положения выключателей (например, \textit{«В включен»} и \textit{«В отключен»}) вводится программным ключом \textbf{«Контр. Р В1»}. Также имеется возможность выбрать количество разъединителей, участвующих в формировании сигналов включенного или отключенного положения выключателей, для выключателя В1 она задается программным переключателем \textbf{«Схема Р В1»}.
\item
Также функция ФПВ пофазно контролирует состояние блок-контактов выключателя и в случае расхождения полюсов формирует сигнал \textit{«Внутр. пуск. ЗНФ»}.
\item
Функциональная схема логики ФПВ В1 приведена на рисунке \ref{pict_fpv:fpv1}. Уставки ФПВ В1 представлены в таблице \ref{app_set:fpv_q1}.
}{}
\ifbool{isDZAT}{
\item
Контроль положения разъединителей при формировании сигналов включенного или отключенного положения выключателей (например, \textit{«В1(ЛВ) ВН включен»} и \textit{«В1(ЛВ) ВН отключен»} для выключателя В1(ЛВ) ВН) вводится программным ключом \textbf{«Контр. Р В1(ЛВ) ВН»}. Также имеется возможность выбрать количество разъединителей, участвующих в формировании сигналов включенного или отключенного положения выключателей, для выключателя В1 она задается программным переключателем \textbf{«Схема Р В1(ЛВ) ВН»}.
\item
Также функция ФПВ пофазно контролирует состояние блок-контактов выключателя и в случае расхождения полюсов формирует сигнал \textit{«Внутр. пуск. ЗНФ В1(ЛВ) ВН»}.
\item
Функциональная схема логики ФПВ В1(ЛВ) ВН приведена на рисунке \ref{pict_fpv:fpv1}. Функциональная схема логики ФПВ В1(ЛВ) СН выполнена аналогично. Уставки ФПВ В1(ЛВ) ВН представлены в таблице \ref{app_set:fpvvn1}. Уставки ФПВ В1(ЛВ) СН представлены в таблице \ref{app_set:fpvsn1}
}{}


\vspace{2mm}
\begin{figure}[H]
\centering
\includegraphics[width=1\textwidth,height=1\textheight,keepaspectratio]{\fbpath/03.01 ЛВ/761.3381 ФПВ12/img/fpv1.pdf}
\captionof{figure}{Функциональная схема логики ФПВ \Qone}\label{pict_fpv:fpv1}
\end{figure}

\ifbool{isLV}{
\item
Функциональная схема логики ФПВ В2 приведена на рисунке \ref{pict_fpv:fpv2}. Описание работы функциональной схемы логики аналогично ФПВ В1. Уставки ФПВ В2 представлены в таблице \ref{app_set:fpv_q2}.
}{}
\ifbool{isDZAT}{
\item
Функциональная схема логики ФПВ В2(ОВ) ВН приведена на рисунке \ref{pict_fpv:fpv2}. Функциональная схема логики ФПВ В1(ЛВ) СН выполнена аналогично. Уставки ФПВ В2(ОВ) ВН представлены в таблице \ref{app_set:fpvvn2}. Уставки ФПВ В2(ОВ) СН представлены в таблице \ref{app_set:fpvsn2}
}{}

\vspace{2mm}
\begin{figure}[H]
\centering
\includegraphics[width=1\textwidth,height=1\textheight,keepaspectratio]{\fbpath/03.01 ЛВ/761.3381 ФПВ12/img/fpv2.pdf}
\captionof{figure}{Функциональная схема логики ФПВ \Qtwo}\label{pict_fpv:fpv2}
\end{figure}

\end{enumerate}